\scanpage{179}
\begin{center}- 29 -\end{center}\hfill IX

démonstration de a) (en se ramenant au cas $S$ local noethérien complet,
et utilisant SGA 3 XV 1.6 c), X 3.2, IX 3.6, IX 5.2), savoir :
\[
 (3.1.4.1)\qquad\operatorname{rang\,mult}({}_nK_s)\mid
 \operatorname{rang\,mult}({}_nK_t)\qquad,
\]
où rang mult désigne le rang du plus grand sous-groupe de type multiplicatif du groupe algébrique commutatif envisagé (SGA 3 XVII 7.2.1) ;
on notera qu'un groupe algébrique de type multiplicatif de torsion est
fini. La formule (3.1.4.1) est valable pour tout entier $n\geq1$. On
conclut, grâce à SGA XVII 4.6.1 (vi), que dans c), si $K_t$ est unipotent
(SGA 3 VII 1.1, il en est de même de $K_s$.

\uline{Proposition} 3.2. \uline{Soient} $S$ \uline{un schéma localement noethérien régulier
intègre de dimension} $1$, $K$ \uline{son corps des fonctions rationnelles}, $A_K$
\uline{un schéma abélien sur} $K$, $A$ \uline{son modèle de Néron sur} $S$, \uline{qui est un
schéma en groupes lisse et séparé sur} $S$ (\uline{et quasi-projectif sur} $S$
\uline{si} $S$ \uline{est noethérien} [ 23 , VIII 2, $3^\circ$]). \uline{Les conditions suivantes
sont équivalentes} :

(i) \uline{Pour tout} $s\in S$, $A_s^{\unsure{o|0}}$ \uline{est de rang unipotent nul}, \uline{i.e.}
(2.1.5) \uline{est une extension d'un schéma abélien par un tore}.

(ii) \uline{Il existe un schéma en groupes} $G$ \uline{lisse}, \uline{séparé de type
fini sur} $S$, \uline{prolongeant} $A_K$, \uline{et dont les fibres sont de rang unipotent nul}.

\uline{De plus}, \uline{si} $G$ \uline{est comme dans} (ii), \uline{alors l'unique morphisme}
$G\longrightarrow A$ \uline{induisant l'isomorphisme donné sur les fibres génériques est
un isomorphisme de} $G$ \uline{sur un sous-groupe ouvert de} $A$, \uline{donc induit un
isomorphisme des composantes neutres}
$G^{\unsure{o|0}}\xrightarrow{\sim}A^{\unsure{o|0}}$.
\margin{347}
\margin{347}
\margin{178}
